\section{Introduction}
\label{sec:s2-introduction}

Distributed formation control depends on information that is both physically
relevant and sufficiently fresh. Periodic communication ignores whether the
plant state has changed, while a conventional event trigger can reset its
local error after transmission even when no neighbor received the packet.
Age-of-information (AoI) scheduling addresses freshness, but the sender does
not directly observe receiver AoI when DATA and acknowledgement (ACK) traffic
experience delay, loss, queueing, and contention. In a broadcast swarm this
problem is receiver specific: one DATA frame may refresh several neighbors,
yet every neighbor can hold a different generation and return feedback through
a different reverse path.

These facts create a coupled communication--control question. When should a
UAV broadcast a new state, and when is a separate ACK worth occupying the same
medium? A standalone ACK carries no plant state. Its only possible benefit is
mediated: earlier confirmation changes sender belief, which can change future
DATA generation, admission, service, or collisions. The same confirmation can
therefore save redundant DATA or suppress a useful refresh. ACK rate and
confirmation delay alone do not determine its control value.

This paper develops and tests a causal semantic-freshness framework for that
setting. Each sender keeps a last-transmitted state for innovation and a
receiver-specific, ACK-confirmed generation time for freshness. Broadcast DATA
is combined with cumulative ACK entries that can be piggybacked or transmitted
separately. The network is an explicit discrete-event shared medium with
multi-slot frames, finite queues, collision/interference and carrier-sense
graphs, retry limits, background occupancy, and independent bursty DATA/ACK
links. Thus a communication claim is made in airtime and contention terms,
not by counting logical directed-link updates as independent transmissions.

The investigation is structured to expose failure rather than tune it away.
The first preregistered shared-medium holdout shows that broadcast aggregation
substantially improves over causal unicast, but the proposed default is
dominated by a delayed-ACK belief point in every primary regime and collapses
at $N=20$ with fixed access probability. Separate development and holdout
blocks isolate analytical access scaling, remove a harmful load guard, and
then examine whether standalone feedback adds value beyond piggyback. Policy-
level shadows and state-cloned branch replay distinguish correlation,
full-policy effects, and the causal effect of one focal ACK decision.

The contributions are:

\begin{enumerate}
\item An integrated receiver-confirmed broadcast protocol that separates innovation from
freshness, aggregates cumulative ACK obligations, and preserves causal
sender knowledge under a finite shared-medium implementation.

\item A mathematical chain from ACK integrity to conservative sender-side
age, semantic state-error envelopes, stochastic age tails, and formation ISS,
together with an exact ACK-airtime identity and an explicit identification
boundary for branch-at-decision replay.

\item A preregistered mechanism study that retains three consequential
negative results: the initial default is not frontier optimal; fixed access
collapses at $N=20$; and positive causal confirmation lead is not a freshness
or control-benefit certificate.

\item A MAC-conditioned feedback hypothesis evaluated by disjoint directional
validation, a matched factorial holdout, and a preregistered operating-envelope
challenge. The base-cell interaction is supported for RMSE and offered airtime,
but only 1/8 envelope contexts passes the full rule; configured MAC type alone
is therefore rejected as a general selector.
\end{enumerate}

The claim is deliberately narrower than universal superiority. The interaction
is validated only in the declared N5 Moderate base cell, while its general
selector extension fails the frozen operating-envelope gate.
The $N=20$ focal-decision
effect is not identified because eligible events are rare. The channel and
energy parameters are simulated rather than measured, and no security,
hardware, HIL, flight, or population safety claim is made.
